\section{Training}
\label{sec:training}

\subsection{SFT Stages}

The released AmphionASR checkpoint is the output of a multi-stage
supervised fine-tuning (SFT) recipe. Each stage progressively
unlocks more of the model and adds more demanding tasks:

\begin{enumerate}
  \item \textbf{Stage~1 -- Audio--Text Alignment.}
        The audio encoder and the LLM are both frozen; only the
        modality projector (Section~\ref{subsec:projector}) and the
        four learnable modality-boundary prompts receive gradients. The
        objective is to map speech representations into the LLM's
        token space without disturbing either the encoder's acoustic
        priors or the LLM's text capabilities.
  \item \textbf{Stage~2 -- Continue SFT.} The projector, the
        boundary prompts and the LLM are jointly fine-tuned on a
        balanced mixture of Mandarin and English transcription; the
        audio encoder remains frozen. The LLM is updated through a
        LoRA adapter (rank $64$, $\alpha=128$, applied to all linear
        projections); the released checkpoint folds the adapter back
        into the base weights for distribution, which is why the
        published weights contain no separate LoRA tensors. SFT is
        driven by ms-swift~\cite{msswift}.
  \item \textbf{Stage~3 -- Capability SFT.} On top of the Stage~2
        checkpoint we continue SFT on the full multi-task mixture
        (Section~\ref{sec:data}): Mandarin, English and multilingual
        ASR, hotword-conditioned ASR with the online distractor and
        hard-negative pool of Section~\ref{subsec:hotword-pool},
        target-speaker ASR (Section~\ref{subsec:tsasr-synth}) and the
        anti-hallucination negatives. Hyper-parameters are inherited
        from Stage~2 unless noted.
\end{enumerate}

After Stages~1--3 the released checkpoint additionally goes through
a degradation-robustness SFT stage trained on Voice-in-the-Wild-2M~\cite{xie2026mega},
a large-scale dataset covering diverse real-world acoustic degradations
for the plain-ASR branch. Unlike the earlier stages, the audio encoder
is \emph{unfrozen} in this stage. The motivation is that real-world
degradations (reverberation, background noise, codec artefacts) corrupt
the acoustic signal before it reaches the encoder; a frozen encoder
produces distorted speech embeddings that carry degradation-induced
artefacts into the LLM, interfering with its transcription capability.
Jointly fine-tuning the encoder on degraded speech allows it to learn
degradation-robust representations, so that the embeddings passed to
the LLM remain semantically clean regardless of recording condition.
Its external evaluation against eight other systems on
Voices-in-the-Wild-Bench~\cite{xie2026mega} is reported in
Section~\ref{sec:noise}.

\subsection{Group Relative Policy Optimisation}
\label{subsec:grpo}

On top of the degradation-robustness SFT stage, we apply a
reinforcement-learning step based on Group Relative Policy
Optimisation (GRPO) \cite{lakomkin2025contextual} with two
rule-based rewards computed directly on the assistant text. The
numbers reported in Sections~\ref{sec:experiments}--\ref{sec:noise}
are produced from the SFT checkpoint described above; a dedicated
SFT vs.\ SFT+GRPO ablation is deferred to a follow-up revision.

\paragraph{Motivation.}
SFT on ASR is teacher-forced, and two failure modes of the resulting
decoder are not directly addressed by more SFT epochs.
\emph{(i)~Boundary-discrimination samples}~-- silent inputs that
should produce an empty hypothesis, and hotword decisions about
whether a candidate term was actually spoken~-- carry a token-count
weight that is small relative to typical transcripts, so the
cross-entropy gradient does not specifically reward correct boundary
behaviour. \emph{(ii)~Metric-train gap}~-- CER is a sequence-level
edit-distance whose minimum-CER hypothesis is not guaranteed to be
the minimum-cross-entropy hypothesis; near the SFT convergence point
this gap is what remains to be optimised. A sequence-level reward
signal addresses both issues directly.

\paragraph{ASR accuracy reward.}
The accuracy reward is the negative CER clipped to $[0, 1]$,
\begin{equation}
  R_\mathrm{asr}(\hat{y}, y) = \max(0,\, 1 - \cer(\hat{y}, y)),
  \label{eq:reward-asr}
\end{equation}
with three boundary cases. When both the reference and the
hypothesis are empty, $R_\mathrm{asr}=1$ (silent input correctly
produces silence). When the reference is empty but the hypothesis
is not, $R_\mathrm{asr}=0$ (hallucination on silence is penalised
maximally). When $\cer(\hat{y}, y) > 1$ -- e.g.\ severe insertion
loops -- the reward is clipped to $0$ rather than allowed to take
arbitrarily large negative values; otherwise a single broken
completion would dominate the within-group $z$-score and flatten
the advantage of the remaining completions. We use CER rather than
WER throughout because our training and test sets are mixed
Mandarin and English; CER is the only character-level metric that
is consistently defined across the two languages and the
code-switching slice.

\paragraph{Hotword match-accuracy reward.}
Given a candidate set $C$ parsed from the prompt's hotword line,
let $\mathrm{Pred}(\hat{y}; C)$ and $\mathrm{Ref}(y; C)$ be the
subsets of $C$ actually present in the hypothesis and the
reference, respectively. Both are extracted with a longest-first
masked match that prevents nested double-counting (so ``Beijing
duck'' is matched as one term, not as ``Beijing'' plus ``Beijing
duck''). The hotword reward is the average over $C$ of the binary
agreement
\begin{equation}
  R_\mathrm{hw}(\hat{y}, y; C) = \frac{1}{|C|}\sum_{c \in C}
    \mathbf{1}\!\left\{\mathbf{1}[c \in \mathrm{Pred}] = \mathbf{1}[c \in \mathrm{Ref}]\right\}.
  \label{eq:reward-hw}
\end{equation}
We use match accuracy rather than $F_1$ because of an explicit
anti-spray argument. Under $F_1$, a degenerate policy that emits
every candidate term unconditionally attains recall $1$ and
precision equal to $\alpha = |R| / |C|$, the fraction of true
hotwords; the resulting $F_1 = 2\alpha / (1 + \alpha)$ remains a
non-trivial reward even when $\alpha$ is small (e.g.\ $\alpha=0.125$
gives $F_1\!\approx\!0.22$). Under match accuracy the same
degenerate policy attains exactly $\alpha$, which is strictly below
the $1.0$ that the correctly discriminating policy attains; the
hacking path is no longer competitive. The same argument couples
the GRPO stage to the retrieve-then-bias evaluation protocol of
Section~\ref{sec:hotwords}: the retrieve stage is already tuned for
high precision (Section~\ref{subsec:retrieve-algorithm}), but its
output still contains near-duplicate distractors; match accuracy
turns each correctly rejected distractor into a positive reward
signal of $1/|C|$, which is the training-time analogue of the
``abstain'' behaviour the retrieve discriminative-gap check
enforces at inference. Finally, when the prompt has no hotword line,
the candidate set is empty and we return a neutral
$R_\mathrm{hw}=1$ so that plain-ASR samples are not implicitly
penalised on a dimension they do not condition on.

\paragraph{Weighted combination.}
The two rewards combine linearly,
\begin{equation}
  R(\hat{y}, y; C) = w_\mathrm{asr}\, R_\mathrm{asr}(\hat{y}, y)
    + w_\mathrm{hw}\, R_\mathrm{hw}(\hat{y}, y; C),
  \label{eq:reward-total}
\end{equation}
with $(w_\mathrm{asr}, w_\mathrm{hw}) = (1.0, 0.3)$. Two
constraints fix this ratio. \emph{Primary axis.} CER is the
deployment-facing metric and must dominate the gradient direction;
$w_\mathrm{asr}\!\gg\!w_\mathrm{hw}$ ensures that improvements on
hotword matching at the cost of transcript accuracy are penalised.
\emph{Signal floor.} The hotword reward must contribute enough
gradient that the policy continues to learn hotword behaviour at
the SFT convergence point; setting $w_\mathrm{hw}=0$ would
collapse GRPO to pure accuracy optimisation. Empirically the
$[0.2, 0.4]$ band satisfies both constraints; this paper takes the
midpoint and does not claim an ablated optimum.

Hyperparameter details are reported in the appendix.
